\section{Two-Element Array Simulation Results}
\label{sec:array-results}

The tuned single-patch element was duplicated to form the \(2\times1\) array.  After duplication, the
array did not retain the exact same resonance as the isolated single element because the two elements
were electromagnetically coupled.  Therefore, the array was tuned separately.  The final array tuning
used \(\texttt{patchL}=36.70~\mathrm{mm}\) and \(\texttt{insetY}=13.35~\mathrm{mm}\), while preserving
\(\texttt{patchW}=47.382~\mathrm{mm}\), \(\texttt{feedW}=4.735~\mathrm{mm}\), and the element spacing
\(d=59.935~\mathrm{mm}\).

\subsection{Array Geometry}
\label{subsec:array-geometry-results}

Figures~\ref{fig:array-top-view-dimensions} and~\ref{fig:array-cross-section-dimensions}
document the final two-element geometry and layer stack. The two elements were arranged along the
CST $x$-axis with center-to-center spacing $d=59.935~\mathrm{mm}$, while both feed lines were aligned
with the $y$-axis. The complete tuned dimensional set is reported in
Table~\ref{tab:cst-array-final-parameters}; therefore, the figures can remain visually clear without a
crowded parameter-list overlay.

\begin{figure}[H]
    \centering
    \maybeincludegraphics[width=1\textwidth]{figures/5g_patch_array/array_2x1/array_2x1_top_view_with_dimensions.png}
    \caption[Top view of the final tuned $2\times1$ array.]{Top view of the final tuned $2\times1$ inset-fed patch array. The final elements used $W=47.382~\mathrm{mm}$, $L=36.70~\mathrm{mm}$, $y_{\mathrm{inset}}=13.35~\mathrm{mm}$, $g_{\mathrm{inset}}=1.000~\mathrm{mm}$, and center-to-center spacing $d=59.935~\mathrm{mm}$.}
    \label{fig:array-top-view-dimensions}
    \label{fig:array-top-view}
\end{figure}

\begin{figure}[H]
    \centering
    \maybeincludegraphics[width=0.96\textwidth]{figures/5g_patch_array/array_2x1/array_2x1_cross_section_with_dimensions.png}
    \caption[Layered cross-sectional view of the final $2\times1$ array.]{Layered cross-sectional perspective of the final $2\times1$ array. Both elements use the same $0.035~\mathrm{mm}$ bottom metallization, $1.524~\mathrm{mm}$ DiClad 880-IM substrate, and $0.035~\mathrm{mm}$ top metallization stack.}
    \label{fig:array-cross-section-dimensions}
\end{figure}

\begin{table}[H]
\centering
\caption{Layer stack used in the final $2\times1$ array model.}
\label{tab:array-layer-stack}
\small
\renewcommand{\arraystretch}{1.16}
\begin{tabular}{@{}lccc@{}}
\toprule
\textbf{Layer} & \textbf{$z_{\min}$ (mm)} & \textbf{$z_{\max}$ (mm)} & \textbf{Thickness (mm)} \\
\midrule
Ground metallization & 0.000 & 0.035 & 0.035 \\
DiClad 880-IM substrate & 0.035 & 1.559 & 1.524 \\
Top metallization & 1.559 & 1.594 & 0.035 \\
\bottomrule
\end{tabular}
\end{table}

\begin{figure}[H]
    \centering
    \maybeincludegraphics[width=0.80\textwidth]{figures/5g_patch_array/array_2x1/array_2x1_port_view.png}
    \caption{Two-port CST model of the $2\times1$ array. Each patch element is fed by a separate waveguide port, enabling both two-port $S$-parameter extraction and independent phase excitation for beam steering.}
    \label{fig:array-port-view}
\end{figure}

\subsection{Array Input Match and Mutual Coupling}

The first array simulation using the single-patch tuned length showed a downward resonant-frequency
shift.  This is a direct consequence of mutual coupling and the changed boundary environment when a
second element is introduced.  Table~\ref{tab:array-tuning-log} summarizes the array tuning steps.  The
final Run~05 represents the best compromise between input matching, resonance near the target
frequency, and mutual coupling.

\begin{table}[H]
\centering
\caption{Two-element array tuning log.}
\label{tab:array-tuning-log}
\footnotesize
\setlength{\tabcolsep}{4pt}
\renewcommand{\arraystretch}{1.25}

\begin{tabularx}{\textwidth}{@{}
>{\raggedright\arraybackslash}p{0.16\textwidth}
>{\raggedright\arraybackslash}p{0.18\textwidth}
>{\raggedright\arraybackslash}p{0.27\textwidth}
>{\centering\arraybackslash}p{0.11\textwidth}
>{\raggedright\arraybackslash}X
@{}}
\toprule
\textbf{Run} &
\textbf{Geometry} &
\textbf{\(S\)-parameter behavior near target} &
\textbf{Coupling} &
\textbf{Observation} \\
\midrule

Array Run 01 &
\(L=38.40~\mathrm{mm}\), \(y_{\mathrm{inset}}=14.11~\mathrm{mm}\) &
Resonance near \(2.399~\mathrm{GHz}\). &
\(\approx -13.7~\mathrm{dB}\) &
The resonance shifted downward after forming the two-element array. \\

Array Run 02 &
\(L=36.85~\mathrm{mm}\), \(y_{\mathrm{inset}}=13.54~\mathrm{mm}\) &
\(S_{11}\approx -16.9~\mathrm{dB}\), \(S_{22}\approx -17.8~\mathrm{dB}\). &
\(\approx -14.7~\mathrm{dB}\) &
The matching improved, but the resonance remained slightly below the target frequency. \\

Array Run 03/04 &
\(L=36.65~\mathrm{mm}\), \(y_{\mathrm{inset}}=13.54/13.20~\mathrm{mm}\) &
Improved frequency alignment; port balance varied. &
\(\approx -15~\mathrm{dB}\) &
These runs were used to study the sensitivity of the array response to the feed inset position. \\

Final Run 05 &
\(L=36.70~\mathrm{mm}\), \(y_{\mathrm{inset}}=13.35~\mathrm{mm}\) &
\(S_{11}\approx -11.9~\mathrm{dB}\), \(S_{22}\approx -15.2~\mathrm{dB}\). &
\(\approx -15~\mathrm{dB}\) &
This was selected as the final tuned array for the beam-steering simulations. \\

\bottomrule
\end{tabularx}
\end{table}

Figures~\ref{fig:array-s11-s12} and~\ref{fig:array-s21-s22} show the final two-port
\(S\)-parameters for Run~05.  Near the operating frequency, the array satisfies the common
\(-10~\mathrm{dB}\) input-match criterion for both ports.  The mutual coupling is approximately
\(-15~\mathrm{dB}\), which is an important array-level result because it explains why the array had to
be retuned separately from the isolated element.

\begin{figure}[H]
    \centering
    \maybeincludegraphics[width=0.96\textwidth]{figures/5g_patch_array/array_2x1/array_2x1_sparameters_tuned_run05_s11_s12.png}
    \caption{Final Run~05 \(S_{11}\) and \(S_{12}\) for the tuned \(2\times1\) array.  The final array achieved \(S_{11}\approx -11.93~\mathrm{dB}\) near \(2.498~\mathrm{GHz}\) and mutual coupling \(S_{12}\approx -15.19~\mathrm{dB}\) near \(2.494~\mathrm{GHz}\).}
    \label{fig:array-s11-s12}
\end{figure}

\begin{figure}[H]
    \centering
    \maybeincludegraphics[width=0.96\textwidth]{figures/5g_patch_array/array_2x1/array_2x1_sparameters_tuned_run05_s21_s22.png}
    \caption{Final Run~05 \(S_{21}\) and \(S_{22}\) for the tuned \(2\times1\) array.  The second port achieved \(S_{22}\approx -15.20~\mathrm{dB}\) near \(2.497~\mathrm{GHz}\), and the coupling from Port~1 to Port~2 was approximately \(S_{21}\approx -14.57~\mathrm{dB}\) near \(2.499~\mathrm{GHz}\).}
    \label{fig:array-s21-s22}
\end{figure}

\subsection{Individual Element Farfield Responses}

Before combining the port excitations for beam steering, CST produced farfield responses for each
individual port excitation.  These responses were used to verify that each port excited a radiating patch
element and to provide the element-level data used by the post-processing combination step.

\begin{figure}[H]
    \centering
    \maybeincludegraphics[width=0.78\textwidth]{figures/5g_patch_array/array_2x1/array_2x1_3d_farfield_excitation1.png}
    \caption{Array farfield response for excitation of Port~1 only.  This result is used as one of the two element responses in the combined beam-steering post-processing.}
    \label{fig:array-ff-excitation1}
\end{figure}

\begin{figure}[H]
    \centering
    \maybeincludegraphics[width=0.78\textwidth]{figures/5g_patch_array/array_2x1/array_2x1_3d_farfield_excitation2.png}
    \caption{Array farfield response for excitation of Port~2 only.  The two individual excitation results are later combined using amplitude and phase weights.}
    \label{fig:array-ff-excitation2}
\end{figure}

\begin{table}[H]
\centering
\caption{Two-element array final simulation summary.}
\label{tab:array-final-summary}
\begin{tabular}{@{}ll@{}}
\toprule
\textbf{Metric} & \textbf{Final CST Value} \\
\midrule
Target frequency & \(2.501~\mathrm{GHz}\) \\
Final array \texttt{patchL} & \(36.70~\mathrm{mm}\) \\
Final array \texttt{insetY} & \(13.35~\mathrm{mm}\) \\
Element spacing, \(d\) & \(59.935~\mathrm{mm}\) \\
\(S_{11}\) near target & \(\approx -11.93~\mathrm{dB}\) near \(2.498~\mathrm{GHz}\) \\
\(S_{22}\) near target & \(\approx -15.20~\mathrm{dB}\) near \(2.497~\mathrm{GHz}\) \\
Mutual coupling, \(S_{12}\) & \(\approx -15.19~\mathrm{dB}\) near \(2.494~\mathrm{GHz}\) \\
Mutual coupling, \(S_{21}\) & \(\approx -14.57~\mathrm{dB}\) near \(2.499~\mathrm{GHz}\) \\
Combined broadside peak directivity & \(\approx 9.655~\mathrm{dBi}\) (CST displayed value) \\
Metal model for tuned run & PEC approximation \\
Dielectric loss & Retained through \(\tan\delta=0.0009\) \\
\bottomrule
\end{tabular}
\end{table}
